\begin{figure}[H]
    \centering
    \begin{tikzpicture}[x=0.78cm, y=1.0cm]

        \lanerow{0}{v}{7,2,5,8,1,6,3,4}{lane}
        \node[rowlab] at (v0.west) {$v$};

        \lanerow{-1.15}{p}{5,8,7,2,3,4,1,6}{lane}
        \node[rowlab] at (p0.west) {$v'$};

        \foreach \a/\b in {0/2,1/3,2/0,3/1,4/6,5/7,6/4,7/5} {
            \draw[flowmute] (v\a.south) -- (p\b.north);
        }

        \lanerow{-2.30}{lo}{5,2,5,2,1,4,1,4}{lanereg}
        \node[rowlab] at (lo0.west) {$\min(v,v')$};

        \lanerow{-3.05}{hi}{7,8,7,8,3,6,3,6}{lanecac}
        \node[rowlab] at (hi0.west) {$\max(v,v')$};

        \lanerow{-3.85}{m}{0,0,1,1,0,0,1,1}{lane}
        \node[rowlab] at (m0.west) {mask};

        \foreach \i/\val in {0/5,1/2,4/1,5/4} { \node[lanereg] (r\i) at (\i,-5.00) {$\val$}; }
        \foreach \i/\val in {2/7,3/8,6/3,7/6} { \node[lanecac] (r\i) at (\i,-5.00) {$\val$}; }
        \node[rowlab] at (r0.west) {result};

        \foreach \i in {0,...,7} { \node[idxlab] at (r\i.south) {\i}; }

        \foreach \i in {0,1,4,5} { \draw[flowmute] (lo\i.south) to[out=-90,in=90] (r\i.north); }
        \foreach \i in {2,3,6,7} { \draw[flowmute] (hi\i.south) to[out=-90,in=90] (r\i.north); }

        \node[notemute, anchor=west] at (8.25,-0.58) {shuffle by $j=2$};
        \node[notemute, anchor=west] at (8.25,-2.68) {min / max};
        \node[notemute, anchor=west] at (8.25,-4.45) {select with mask};

        \draw[brace] (8.15,-2.05) -- (8.15,-3.30);
        \draw[brace] (8.15,-3.60) -- (8.15,-5.25);
        \draw[brace] (8.15,0.25) -- (8.15,-1.40);

    \end{tikzpicture}
    \caption[The fundamental compare-exchange motif]{The fundamental ``shuffle, min, max, select''
    motif for compare-exchange at distance $j=2$, in an eight-lane register. The
    contents are permuted so that each element ends up opposite its counterpart,
    minimum and maximum are computed simultaneously for all lanes, and the result
    is reassembled by selecting the minimum where the corresponding mask bit
    is zero and the maximum where it is one. No branch intervenes and no value
    is examined individually.}
    \label{fig:cpu-simd-motif}
\end{figure}
